\documentclass[aspectratio=169,11pt]{beamer}
% Shared CMPSC 480 Beamer preamble.
\usepackage[T1]{fontenc}
\usepackage[utf8]{inputenc}
\usepackage{lmodern}
\usepackage{amsmath,amssymb,mathtools}
\usepackage{booktabs,array,tabularx}
\usepackage{xcolor}
\usepackage{listings}
\usepackage{tikz}
\usetikzlibrary{arrows.meta,positioning,shapes.geometric}

% Ignore only negligible TeX box diagnostics that are below visual tolerance.
% Larger layout defects still appear in the build log and in rendered QA.
\vfuzz=3pt
\hbadness=2000

\definecolor{courseink}{HTML}{111111}
\definecolor{coursegray}{HTML}{EDEDED}
\definecolor{courserule}{HTML}{B8BCC4}
\definecolor{courseblue}{HTML}{3D8DFF}
\definecolor{courselightblue}{HTML}{D0EDFA}
\definecolor{coursemuted}{HTML}{5C6470}
\definecolor{coursegreen}{HTML}{0B7A53}
\definecolor{coursered}{HTML}{B3261E}

\usetheme{default}
\usefonttheme{professionalfonts}
\setbeamertemplate{navigation symbols}{}
\setbeamersize{text margin left=0.65cm,text margin right=0.65cm}

\setbeamercolor{normal text}{fg=courseink,bg=white}
\setbeamercolor{frametitle}{fg=courseink,bg=white}
\setbeamercolor{title}{fg=courseink,bg=white}
\setbeamercolor{subtitle}{fg=coursemuted,bg=white}
\setbeamercolor{structure}{fg=courseblue}
\setbeamercolor{alerted text}{fg=coursered}
\setbeamercolor{block title}{fg=courseink,bg=courselightblue}
\setbeamercolor{block body}{fg=courseink,bg=coursegray}
\setbeamercolor{example text}{fg=coursegreen}

\setbeamerfont{title}{size=\fontsize{25}{29}\selectfont,series=\bfseries}
\setbeamerfont{subtitle}{size=\fontsize{13}{16}\selectfont}
\setbeamerfont{frametitle}{size=\fontsize{19}{22}\selectfont,series=\bfseries}
\setbeamerfont{normal text}{size=\fontsize{11}{14}\selectfont}
\setbeamerfont{footline}{size=\fontsize{7.5}{9}\selectfont}

\setbeamertemplate{frametitle}{%
  \nointerlineskip
  % Use content-driven height so a long assessment title can wrap without
  % being clipped by a fixed-height title bar.
  \begin{beamercolorbox}[wd=\paperwidth,sep=0.17cm,leftskip=0.48cm,rightskip=0.48cm]{frametitle}
    \insertframetitle
  \end{beamercolorbox}
  \vspace{-0.08cm}
  {\color{courserule}\rule{\textwidth}{0.35pt}}
}

\setbeamertemplate{footline}{%
  \leavevmode
  \begin{beamercolorbox}[wd=\paperwidth,ht=2.6ex,dp=1.1ex,leftskip=.65cm,rightskip=.65cm]{author in head/foot}
    \color{coursemuted}\insertshorttitle\hfill\insertframenumber/\inserttotalframenumber
  \end{beamercolorbox}
}

\setbeamertemplate{itemize item}{\color{courseblue}\small$\blacktriangleright$}
\setbeamertemplate{itemize subitem}{\color{courseblue}\scriptsize$\bullet$}
\setbeamertemplate{enumerate item}{\color{courseblue}\insertenumlabel.}

\lstdefinestyle{coursepython}{
  language=Python,
  basicstyle=\ttfamily\fontsize{8.5}{10}\selectfont,
  keywordstyle=\color{courseblue}\bfseries,
  commentstyle=\color{coursemuted}\itshape,
  stringstyle=\color{coursegreen},
  showstringspaces=false,
  columns=fullflexible,
  keepspaces=true,
  frame=single,
  rulecolor=\color{courserule},
  backgroundcolor=\color{coursegray},
  breaklines=true,
  tabsize=4,
  xleftmargin=0.15cm,
  xrightmargin=0.15cm,
  framexleftmargin=0.15cm,
  framexrightmargin=0.15cm
}
\lstset{style=coursepython}

\newcommand{\points}[1]{\hfill{\color{courseblue}\textbf{[#1 points]}}}
\newcommand{\deliverable}[1]{\textbf{\color{courseblue}#1}}
\newcommand{\code}[1]{\texttt{#1}}
\newcommand{\keyidea}[1]{%
  \begin{beamercolorbox}[rounded=false,sep=0.55em,wd=\linewidth]{block body}
    \textbf{Key idea.} #1
  \end{beamercolorbox}
}
\newcommand{\answerlabel}{\textbf{\color{coursegreen}Answer.}\ }
\newcommand{\gradinglabel}{\textbf{\color{courseblue}Grading guidance.}\ }

\newenvironment{questionblock}[1]
  {\begin{block}{#1}}
  {\end{block}}

\newenvironment{solutionblock}[1]
  {\begin{block}{#1}}
  {\end{block}}

\AtBeginSection[]{
  \begin{frame}
    \vfill
    {\usebeamerfont{title}\color{courseink}\insertsectionhead\par}
    \vspace{0.4cm}
    {\color{courseblue}\rule{0.32\paperwidth}{3pt}}
    \vfill
  \end{frame}
}


% CMPSC480_TOTAL_POINTS: 100
% CMPSC480_ITEM_IDS: 1,2,3,4,5,6,7
\title[Midterm Questions]{CMPSC 480 Midterm Examination}
\subtitle{Cumulative assessment of Weeks 1-7}
\author{CMPSC 480 - Student Question Deck}
\date{}

\begin{document}


\begin{frame}
  \titlepage
\end{frame}

\begin{frame}{Assessment overview}
\begin{columns}[T,onlytextwidth]
\begin{column}{0.62\textwidth}
\Large Solve applied, mathematical, algorithmic, and engineering problems under explicit assumptions.

\vspace{0.35cm}
\normalsize
Coverage: Modules A and B: agents, search, adversarial reasoning, MDPs, RL, approximation, and Bayesian inference
\end{column}
\begin{column}{0.33\textwidth}
\begin{block}{Points}100 total\end{block}
\begin{block}{Expected time}110 minutes\end{block}
\begin{block}{Submission}Individual examination PDF\end{block}
\end{column}
\end{columns}
\end{frame}

\begin{frame}{Learning objectives}
\begin{itemize}
  \item Formulate agent and search problems precisely.
  \item Trace cost-sensitive and adversarial algorithms.
  \item Apply Bellman and temporal-difference updates.
  \item Reason about function approximation stability.
  \item Compute Bayesian posteriors and factor operations.
  \item Integrate methods while preserving assumptions.
\end{itemize}
\end{frame}

\begin{frame}{Requirements checklist}
\small
\begin{itemize}
  \item Work independently and show intermediate algebra, frontier state, or factor state.
  \item One double-sided handwritten letter-size reference sheet is permitted.
  \item A nonprogrammable calculator is permitted; networked devices and communication are prohibited.
  \item State every tie-breaking, terminal, reward, and zero-division convention you use.
  \item Answers receive credit for justified reasoning; an unsupported final number may receive little credit.
  \item Academic integrity rules of the course and institution apply throughout the examination.
\end{itemize}
\end{frame}

\begin{frame}{Task 1 - Agent and search formulation (10 points)}
\small
A delivery robot operates in a partially mapped building with battery limits and moving people.

\vspace{0.2cm}
\begin{enumerate}
\item[\textbf{a.}] Give one concrete element for each PEAS component. \hfill{\color{courseblue}\textbf{[3 points]}}
\item[\textbf{b.}] Define states, actions, transition model, goal test, and path cost for deterministic route planning. \hfill{\color{courseblue}\textbf{[3 points]}}
\item[\textbf{c.}] Explain why rational behavior does not imply omniscience and name two defensive interface checks. \hfill{\color{courseblue}\textbf{[4 points]}}
\end{enumerate}
\end{frame}

\begin{frame}{Task 2 - UCS and A-star (15 points)}
\small
Graph edges are S-A:4, S-B:1, B-A:1, A-G:3, and B-G:8. Heuristic values are h(S)=5, h(A)=3, h(B)=4, h(G)=0.

\vspace{0.2cm}
\begin{enumerate}
\item[\textbf{a.}] Trace UCS until a goal is removed, including cheaper rediscovery. \hfill{\color{courseblue}\textbf{[5 points]}}
\item[\textbf{b.}] Trace A-star priorities and return the selected path. \hfill{\color{courseblue}\textbf{[5 points]}}
\item[\textbf{c.}] Is the heuristic admissible and consistent? Justify from the graph. \hfill{\color{courseblue}\textbf{[5 points]}}
\end{enumerate}
\end{frame}

\begin{frame}{Task 3 - Minimax and alpha-beta (15 points)}
\small
A MAX root has MIN children A, B, and C. Their leaves in search order are A:[6,4], B:[3,8], C:[5,2].

\vspace{0.2cm}
\begin{enumerate}
\item[\textbf{a.}] Compute every minimax value and the selected action. \hfill{\color{courseblue}\textbf{[5 points]}}
\item[\textbf{b.}] Trace alpha-beta left to right and identify pruned leaves. \hfill{\color{courseblue}\textbf{[5 points]}}
\item[\textbf{c.}] Design a depth-cutoff evaluation with three features and state two tests. \hfill{\color{courseblue}\textbf{[5 points]}}
\end{enumerate}
\end{frame}

\begin{frame}{Task 4 - MDP planning (15 points)}
\small
At state s, action a reaches terminal reward 5 with probability 0.8 and state t with reward 0 and probability 0.2. Let gamma=0.9 and V(t)=2. Action b deterministically yields reward 3 and terminates.

\vspace{0.2cm}
\begin{enumerate}
\item[\textbf{a.}] Compute Q(s,a) under reward-on-transition semantics. \hfill{\color{courseblue}\textbf{[5 points]}}
\item[\textbf{b.}] Compute Q(s,b) and the greedy action. \hfill{\color{courseblue}\textbf{[5 points]}}
\item[\textbf{c.}] Explain why value iteration converges for a finite discounted MDP. \hfill{\color{courseblue}\textbf{[5 points]}}
\end{enumerate}
\end{frame}

\begin{frame}{Task 5 - Q-learning and approximation (15 points)}
\small
A linear agent has w=[0.5,-0.5], f(s,a)=[1,2], reward 2, gamma=0.5, max next Q=3, and alpha=0.2.

\vspace{0.2cm}
\begin{enumerate}
\item[\textbf{a.}] Compute current Q, target, TD error, and new weights. \hfill{\color{courseblue}\textbf{[5 points]}}
\item[\textbf{b.}] Recompute the update if the transition is terminal. \hfill{\color{courseblue}\textbf{[5 points]}}
\item[\textbf{c.}] Name the deadly triad and two observable instability signals. \hfill{\color{courseblue}\textbf{[5 points]}}
\end{enumerate}
\end{frame}

\begin{frame}{Task 6 - Bayesian networks and inference (15 points)}
\small
R is a binary root with P(R)=0.2. S depends on R with P(S|R)=0.1 and P(S|not R)=0.5. W depends on both with P(W|R,S)=0.99 and P(W|R,not S)=0.8.

\vspace{0.2cm}
\begin{enumerate}
\item[\textbf{a.}] Compute P(R,S,W) for all three true. \hfill{\color{courseblue}\textbf{[5 points]}}
\item[\textbf{b.}] Write an exact expression for P(R|W) without evaluating missing CPT rows. \hfill{\color{courseblue}\textbf{[5 points]}}
\item[\textbf{c.}] If factors f(R,S) and g(S,W) are multiplied and S is summed out, what is the result scope and why can order matter? \hfill{\color{courseblue}\textbf{[5 points]}}
\end{enumerate}
\end{frame}

\begin{frame}{Task 7 - Integrated engineering scenario (15 points)}
\small
A robot uses a Bayesian belief over hazards as features for a linear Q-learning policy.

\vspace{0.2cm}
\begin{enumerate}
\item[\textbf{a.}] Specify the belief-to-policy interface and three invariants. \hfill{\color{courseblue}\textbf{[5 points]}}
\item[\textbf{b.}] Design a controlled ablation showing whether belief tracking helps. \hfill{\color{courseblue}\textbf{[5 points]}}
\item[\textbf{c.}] State a response to impossible evidence and a posttraining safety gate. \hfill{\color{courseblue}\textbf{[5 points]}}
\end{enumerate}
\end{frame}

\begin{frame}{Edge cases and defensive programming}
\small
\begin{itemize}
  \item A start state is already a goal.
  \item All minimax root actions tie.
  \item A terminal RL transition has a large next Q.
  \item Bayesian evidence has zero probability.
  \item A factor order creates an unexpectedly large table.
\end{itemize}
\end{frame}

\begin{frame}{Submission instructions}
\small
\begin{itemize}
  \item Write your name and identify every requested problem and subpart.
  \item Show enough intermediate work that another reader can audit the result.
  \item Submit the complete question-response document when time is called.
\end{itemize}
\vspace{0.2cm}
\alert{Do not submit credentials, generated caches, or unapproved libraries.}
\end{frame}

\begin{frame}{Grading rubric - part 1}
\scriptsize
\begin{tabularx}{\textwidth}{>{\bfseries}p{0.28\textwidth} p{0.08\textwidth} X}
\toprule
Category & Points & Required evidence \\
\midrule
Agent and search formulation & 10 & PEAS analysis, search model, and engineering contract. \\
UCS and A-star & 15 & Priority trace, optimal path, and heuristic argument. \\
Minimax and alpha-beta & 15 & Minimax backups, pruning trace, and evaluation design. \\
MDP planning & 15 & Bellman values, policy choice, and convergence reasoning. \\
\bottomrule
\end{tabularx}
\end{frame}

\begin{frame}{Grading rubric - part 2}
\scriptsize
\begin{tabularx}{\textwidth}{>{\bfseries}p{0.28\textwidth} p{0.08\textwidth} X}
\toprule
Category & Points & Required evidence \\
\midrule
Q-learning and approximation & 15 & Semi-gradient update, terminal change, and stability analysis. \\
Bayesian networks and inference & 15 & Joint probability, posterior structure, and factor operation. \\
Integrated engineering scenario & 15 & Interface, ablation, and safety verification plan. \\
\bottomrule
\end{tabularx}
\end{frame}

\begin{frame}{Reflection prompt}
In one or two sentences:

Identify one assumption that changed the answer to a problem.
\end{frame}

\end{document}
